\section{Design overview}

\sable{} uses three independent secrets:
\[
    t\in\F_q^n,\qquad s\in\F_q^n,
    \qquad r\in\F_q^{n_c}.
\]
The secret $t$ is used for compact input encryption.  The secret $s$ is used for GSW-style matrix ciphertexts and nonlinear evaluation.  The secret $r$ is used only for the code/LPN compaction layer.

The key chain is
\[
    t \longrightarrow s \longrightarrow r.
\]
An expansion key encrypts the coordinates of $\tilde{t}$ under the GSW-style secret $s$.  This turns compact sparse-LPN input ciphertexts into matrix ciphertexts.  A compaction key encrypts the coordinates of $\tilde{s}$ under the code/LPN compaction secret $r$.  This lets the evaluator homomorphically compute the final linear decryption form and return a compact output.

\subsection{Three ciphertext forms}

\paragraph{Input ciphertexts.}
A compact input ciphertext under $t$ has the form
\[
    c=(a,b)=(a,a\transpose t+\mu+e)\in\F_q^{n+1},
\]
where $a$ is $k$-sparse.  It satisfies
\[
    c\cdot\tilde{t}=\mu+e.
\]
This is small and suitable for data-owner encryption.

\paragraph{Evaluation ciphertexts.}
A GSW-style evaluation ciphertext under $s$ is a matrix $C\in\F_q^{N\times N}$, where $N=n+1$, satisfying
\[
    C\tilde{s}=\mu\tilde{s}+e.
\]
Homomorphic addition and multiplication are simply
\[
    C_1+C_2,
    \qquad
    C_1C_2.
\]
The price is that row support and error grow with circuit depth.

\paragraph{Compact output ciphertexts.}
After evaluation, if
\[
    C_*\tilde{s}=f(\mu_1,\ldots,\mu_m)\tilde{s}+e_*,
\]
then the last row $\rho_*$ satisfies
\[
    \rho_*\cdot\tilde{s}=f(\mu_1,\ldots,\mu_m)+e_{*,N}.
\]
The evaluator uses code/LPN encryptions of $\tilde{s}_i$ to homomorphically compute an encryption of $\rho_*\cdot\tilde{s}$.  This is the compaction step.

\subsection{Why the construction is code-based}

Both hardness assumptions are coding-theoretic.
\begin{itemize}
    \item Sparse LPN is a noisy sparse linear-code decoding problem.  It supplies compact input encryption and the GSW-style nonlinear layer.
    \item q-ary LPN/code encryption is a noisy random linear-code primitive.  It supplies linearly homomorphic compaction.
\end{itemize}
No lattice, number-theoretic, or pairing assumption appears in the construction.  This is the intended novelty.

\subsection{Main technical invariant}

The key invariant is that evaluation ciphertexts remain \emph{approximately eigenvector matrices}.  A ciphertext $C$ encrypting $\mu$ satisfies
\[
    C\tilde{s}=\mu\tilde{s}+e.
\]
Then
\[
    (C_1+C_2)\tilde{s}=(\mu_1+\mu_2)\tilde{s}+(e_1+e_2),
\]
and
\[
    C_1C_2\tilde{s}
    =C_1(\mu_2\tilde{s}+e_2)
    =\mu_1\mu_2\tilde{s}+\mu_2e_1+C_1e_2.
\]
Thus multiplication is valid if $C_1e_2$ remains sparse/noisy enough.  This is why row support matters.
